\documentclass[tikz,border=10pt]{standalone}
\usepackage{ctex}
\usepackage{tikz}
\usetikzlibrary{arrows.meta, calc, positioning}

\definecolor{startred}{RGB}{231,76,60}
\definecolor{nodeblue}{RGB}{52,152,219}
\definecolor{arrowgray}{RGB}{102,102,102}

\begin{document}
\begin{tikzpicture}[font=\large, line cap=round, line join=round]
  % Title
  \node at (4,3.2) {\textbf{C题：有向图与DFS遍历}};

  % Nodes
  \node[circle, fill=startred, text=white, minimum size=1.2cm, inner sep=0pt] (n1) at (1.2,2) {1};
  \node[circle, fill=nodeblue, text=white, minimum size=1.2cm, inner sep=0pt] (n2) at (3.5,2.4) {2};
  \node[circle, fill=nodeblue, text=white, minimum size=1.2cm, inner sep=0pt] (n3) at (6,2) {3};
  \node[circle, fill=nodeblue, text=white, minimum size=1.2cm, inner sep=0pt] (n4) at (3.5,0.6) {4};
  \node[circle, fill=nodeblue, text=white, minimum size=1.2cm, inner sep=0pt] (n5) at (6,0.5) {5};

  % Edges with arrows
  \draw[-{Stealth[length=3mm]}, thick, arrowgray] (n1) -- (n2);
  \draw[-{Stealth[length=3mm]}, thick, arrowgray] (n2) -- (n3);
  \draw[-{Stealth[length=3mm]}, thick, arrowgray] (n2) -- (n4);
  \draw[-{Stealth[length=3mm]}, thick, arrowgray] (n4) -- (n3);
  \draw[-{Stealth[length=3mm]}, thick, arrowgray] (n5) -- (n2);

  % Labels
  \node[below=0.15cm of n1, startred, font=\small] {起点(物品1)};
  \node[below=0.1cm of n4, font=\small] {可达物品};

  % Bottom note
  \node at (4,-0.8) {\small DFS从节点1出发，能到达的节点即为可获得物品};
  \node at (4,-1.4) {\small ans = 可达节点数（包含起点）};
\end{tikzpicture}
\end{document}
